
This work demonstrates that dose-response analysis can be applied to LLM curation parameters, treating them as continuous variables rather than discrete configuration choices.

\subsubsection{7.1 Summary}

The main findings at proof-of-concept scale are:

\begin{enumerate}
\item \textbf{Non-monotonic dose-response exists.} The relationship between perplexity thresholds and benchmark performance is cubic ($R^2$ = 0.985), not linear, with estimated peak near p50.
\end{enumerate}

\begin{enumerate}
\item \textbf{Noise dilution explains the left side of the curve.} Unfiltered training shows 40\% higher convergence AUC (p = 0.044, Cohen's d = -0.62).
\end{enumerate}

\begin{enumerate}
\item \textbf{Scale transfer is feasible.} Optimal thresholds transfer from 125M to 1B scale with ratio 0.85.
\end{enumerate}

\begin{enumerate}
\item \textbf{CPDR-optimized configuration improves over defaults.} 1.32\% improvement at PoC scale (mock evaluation).
\end{enumerate}

\subsubsection{7.2 Future Directions}

\textbf{Full-scale replication.} The primary need is full-scale experiments (10B tokens) with actual benchmark evaluation to confirm effect magnitudes.

\textbf{Architecture generalization.} Replication with Llama-style architectures to test whether non-monotonic patterns are architecture-specific.

\textbf{Alternative quality signals.} Testing classifier-based quality scores, BERT perplexity, or semantic filtering.

\textbf{Joint optimization.} Studying perplexity $\times$ deduplication interaction effects.

\subsubsection{7.3 Closing}

Data curation has been treated as a pipeline preprocessing step---important but unoptimized. These findings suggest it deserves systematic attention. The optimal perplexity threshold of approximately p50 provides a starting point, but the dose-response methodology offers a framework for continued optimization as corpora and architectures evolve.
